\subsection{Foundation models in robotics}
Recent advancements in End-to-End Vision-Language-Action (VLA) models adopt a co-training paradigm that integrates a pre-trained VLM backbone with a dedicated action expert~\cite{Black20240AV, Intelligence202505AV}. By harnessing the VLM's extensive world knowledge~\cite{Bai2025Qwen3VLTR} to interpret complex scenes, this approach effectively transforms the backbone into a generalist agent capable of mapping raw observations directly to continuous actions. 
Alternatively, the hierarchical paradigm decouples reasoning from execution~\cite{Shi2025HiRO, Li2025HAMSTERHA, Shentu2024FromLT}. In this framework, a high-level policy, typically a VLM, processes visual observations and task descriptions to generate intermediate goals, such as natural language sub-tasks or latent embeddings. These intermediate representations then condition a low-level policy to output continuous actions. Building upon this hierarchical structure, our work introduces a novel hierarchical memory architecture designed to empower the high-level VLM to generate more precise and context-aware instructions for the low-level controller.

% \subsection{Robotic memory}
% Memory capabilities are fundamental for VLA to successfully execute complex, long-horizon manipulation tasks. 
% Consequently, extensive research has been dedicated to endowing vla with historical context, though most existing implementations remain predominantly image-based~\cite{Huang2022VisualLM}. 
% These works typically follow two paradigms: utilizing external memory banks~\cite{Li2025MAPVLAMP, Torne2025LearningLD, Fang2025SAM2ActIV, Huang2023OutOS, Zhu2020VisionDialogNB, Zheng2024TraceVLAVT, Hu20253DLLMMemLS}, that employ FIFO queues~\cite{Sridhar2025MemERSU} or fused buffers~\cite{Shi2025MemoryVLAPM} to store raw observations, or exploring history compression to distill information into compact representations. Notably, MemER~\cite{Sridhar2025MemERSU} recently advanced the latter by proposing keyframe selection algorithms to better preserve critical information within subtasks.
% Distinct from these methods, we propose a cross-modal, explicitly managed memory bank designed to enhance task consistency and effectively retain long-term temporal context.

\subsection{Robotic memory}
Memory is important for robotic agents operating in long-horizon, partially observed manipulation tasks, where successful execution often depends on retaining past observations, inferred world states, and user-specific context. Prior work has therefore explored a range of mechanisms for incorporating historical information into robotic control.

One major direction introduces memory at the policy level, where historical context is used primarily to improve local action generation, spatial grounding, and short-range temporal consistency. Typical approaches either augment Vision-Language-Action models with retrieval-style memory banks that store and retrieve relevant past context~\cite{Li2025MAPVLAMP, Fang2025SAM2ActIV, Zhu2020VisionDialogNB, Shi2025MemoryVLAPM}, or encode history into more compact forms, such as past tokens or visual traces, for history-conditioned action prediction~\cite{Torne2025LearningLD, Zheng2024TraceVLAVT}. 

Another direction studies memory in more structured, higher-level decision processes. Previous hierarchical embodied-agent systems~\cite{Wang2023JARVIS1OM, Li2024Optimus1HM, Wang2023DescribeEP} have highlighted the value of explicit planning, reflection, and multimodal memory for long-horizon decision making. MemER~\cite{Sridhar2025MemERSU} organizes experience through a FIFO queue together with keyframe selection. These works suggest that memory is useful not only for short-range control, but also for maintaining task-level context over extended horizons.

Our work follows this latter direction, and focuses on a multimodal memory design with explicit memory management.